%% kapa: exam=1 round=32 year=2021 period=1 form=A subject=economics subject_ko=경제학원론 seq=17 number=57
%% source: Q-Net 감정평가사 기출문제 (cst00309) artlSeq=5208397 file=extracted/1008/5208397/A형/2021년도 제32회 감정평가사 1차 1교시 A형.pdf
%% conversion: kordoc 4.24.51 -> tools/gen1.py (auto); review=proofread-sonnet; flags=none
\begin{question}{57}
오염물질을 배출하는 기업 甲과 乙의 오염저감비용은 각각 $TAC_1=200+4X_1^2$, $TAC_2=200+X_2^2$이다. 정부가 두 기업의 총오염배출량을 80톤 감축하기로 결정할 경우, 두 기업의 오염저감비용의 합계를 최소화하는 甲과 乙의 오염감축량은? (단, $X_1$, $X_2$는 각각 甲과 乙의 오염감축량이다.)
\begin{choices}
  \item $X_1 = 8$, $X_2 = 52$
  \item $X_1 = 16$, $X_2 = 64$
  \item $X_1 = 24$, $X_2 = 46$
  \item $X_1 = 32$, $X_2 = 48$
  \item $X_1 = 64$, $X_2 = 16$
\end{choices}
\end{question}
\answer{2}
